% Short running form; the full title is too wide for the head's right box.
\runninghead{The Invitation Carried to the Highways}
\properlane{call-to-the-nations}{The Invitation Refused and Carried to the Highways}
\label{sec:call-to-the-nations}

Heard from its beginning, the Gospel is the story of an invitation: who sent
it, to whom, how often, what the invited did with it, and where it went when
they would not come. The Mass surrounds that story with the Church's memory of
God's dealings with his people. The Lord who speaks in the Introit, who
promises to hear whoever cries to him and says \latin{Attendite, popule meus,
legem meam}, is the king of the parable, who called the first invited through his servants,
the prophets and after them the apostles, and who, when the invited made light
of the call and killed his messengers, sent his servants to the highways. The
Alleluia gives that last command in the psalm's own words, \latin{annuntiate
inter gentes opera eius}, declare his deeds among the Gentiles. The assembly
that hears this Mass is the hall filled from the highways, and the call comes to
it as a warning taken from the history it has just been told. St John
Chrysostom, St Irenaeus and St Hilary of Poitiers each read the parable as the
call passing from the first invited to the nations, with St Jerome beside them;
St Augustine reads the highways in a form of his own; and St Gregory the Great
reads them otherwise.

\subsection{The sequel to the vineyard: St John Chrysostom}

St John Chrysostom reads the parable where St Matthew places it, as the answer
that follows the vineyard. At the end of the parable of the husbandmen Christ
had said that the kingdom ``shall be given to a nation bringing forth the
fruits thereof'', and the marriage feast, Chrysostom says, ``declares next to
what kind of nation''. It ``proclaims beforehand both the casting out of the
Jews, and the calling of the Gentiles; and it indicates together with this also
the strictness of the life required'' (\work{Homilies on Matthew} 69).

His reading is historical and pastoral at once. The difference between the two
parables is the difference between the times. ``For there He appears before
His crucifixion bidding them; but here even after He is slain, He still urges
them, striving to win them over.'' The invitation had come many times. ``And
when were they bidden? By all the prophets; by John again \dots\ by the Son
Himself again'', and after the Ascension not only by words but by actions,
``by Peter, and those with him''. The excuses of farm and merchandise are
meant to sound reasonable, and that is their lesson: ``though the things which
hinder us be necessary, to set the things spiritual at a higher price than
all''. The burning of the city looks ahead to ``the things that took place
under Vespasian and Titus'', and Chrysostom adds that the capture came ``not
straightway after Christ was slain \dots\ but after forty years, that He might
show His long suffering, when they had slain Stephen, when they had put James
to death, when they had spitefully entreated the apostles''.

Chrysostom is careful about the order of the call. Christ sent his apostles to
the Jews first, both before the crucifixion and after it, and they went to the
Gentiles only when they had been driven out; ``hear Paul interpreting this
parable, and saying thus, `It was necessary that the word of God should first
have been spoken to you, but since ye judge yourselves unworthy, lo, we turn to
the Gentiles'\,'' (Acts 13:46). Then, having told the story of the refusal, he
turns it on those who have come in: ``in order that not even these should put
confidence in their faith alone, He discourses unto them also concerning the
judgment to be passed upon wicked actions'', and the garment ``is life and
practice''. The call that passed from the first invited to the highways is not
a privilege for those it reached. It carries the same judgement with it.

\subsection{One King who calls and judges: St Irenaeus}

St Irenaeus used the parable in the second century against teachers who
divided the God of the prophets from the Father of Jesus Christ. For him the
parable shows ``that there is one King and Lord, the Father of all'', and
``that He had from the beginning prepared the marriage for His Son, and used,
with the utmost kindness, to call, by the instrumentality of His servants, the
men of the former dispensation to the wedding feast''. When they would not come
he sent others, and they killed the messengers; he destroyed them and burned
their city, ``but He called together from all the highways, that is, from all
nations, [guests] to the marriage feast of His Son''. From this Irenaeus draws
his conclusion: ``The Lord, therefore, who has called us everywhere by the
apostles, is He who called those of old by the prophets''
(\work{Against Heresies} IV.36.5).

The same argument governs his reading of the parable's end. The God who calls
the unworthy is the God who examines those he has called, ``because nothing
unbecoming or evil pleases Him'', and the judgement falls on the new people as
it fell on the old: ``For as in the former covenant, `with many of them was He
not well pleased;' so also is it the case here, that `many are called, but few
chosen'\,'' (IV.36.6, citing 1 Cor 10:5). Irenaeus alone among the witnesses
makes the parable a proof that the covenant with Israel and the call of the
nations are the work of one God. That proof is what lets the Introit's Lord,
who speaks to ``my people'' in the psalm of Israel's history, be heard as the
same Lord who calls the Gentiles in the Alleluia and the Gospel.

\subsection{The nations at the end of the way: St Hilary}

St Hilary of Poitiers reads the parable in two movements, and the second lifts
the first. The wedding itself is \latin{vitae coelestis et in resurrectione
suscipiendae aeternae gloriae sacramentum}, the mystery of the heavenly life and
of the eternal glory to be received in the resurrection (\work{Commentarius in
Matthaeum} 22, 3). Then he names the persons. \latin{Invitati antea populus
Israel est: in gloriam enim aeternitatis per legem est advocatus}: the first
invited are the people of Israel, called through the law to the glory of
eternity. The servants sent to call them are the apostles, and those sent a
second time \latin{cum praeceptorum conditione} are apostolic men, the
apostles' successors. The fatted bulls are the glorious company of the martyrs,
and the fatlings are spiritual men; and when all these are ready and gathered
into the number of the multitude pleasing to God, \latin{regni coelestis gloria
tamquam nuptiae nuntiantur}, the glory of the heavenly kingdom is announced as
a wedding (22, 4). The farm and the merchandise are the ambition of the world
and the love of money, and the armies that burn the city are the armies of
heaven, which will burn the murderers' whole assembly with eternal fire
(22, 5).

The second movement begins with the highways. \latin{Quoniam vero de judicii
tempore et resurrectionis ista loquitur, sermonem eumdem ad congregationem
gentium retulit}: because the parable speaks of the time of judgement and
resurrection, the king turns the same word to the gathering of the nations.
The first invited were found unworthy, and the servants are sent \latin{ad
exitus viarum}, and Hilary reads the phrase twice over. The way, he says, is
often the time of this world, and its outlets are its end; the servants are
sent to the ends of the roads \latin{quia omnibus retroacta donantur}, because
at that end everything past is forgiven to all, by the gift of grace. The call
then goes out \latin{sine aliqua exceptione}, without any exception, and the
bad come with the good. \latin{Vocatio quidem bonos efficere debuerat, quia
sancta est}: the call ought to have made them good, because it is holy;
\latin{sed per vitium inemendatae voluntatis discrimen est vocatorum}, but by
the fault of an unamended will there is a difference among the called
(22, 6). And he ends the parable with a sentence that turns its last saying:
\latin{Non est igitur paucitas in invitatis, sed raritas in electis}, the few
are not few among the invited but rare among the chosen (22, 7).

\subsection{Peoples and teachings: St Jerome and St Augustine}

St Jerome agrees with the three and adds his own details. The king who made the
marriage is almighty God, and \latin{Facit autem nuptias Domino nostro Jesu
Christo et Ecclesiae, quae tam ex Judaeis, quam ex gentibus, congregata est}:
he makes the wedding for our Lord Jesus Christ and the Church, which is
gathered as much from the Jews as from the Gentiles. The first servant is
Moses, or, if one reads the plural, the prophets; the second sending is the
prophets or the apostles. Those who went off to farm and merchandise are
guilty of a lesser crime than those who turned humanity into cruelty. Jerome
notices that the parable calls the king \latin{homo}, a man, while he invites
and does works of clemency, and only \latin{rex} when he comes to punish. His
armies are either avenging angels, \latin{de quibus in Psalmis scribitur:
Immissiones per angelos pessimos}, as it is written in the Psalms, ``which he
sent by evil angels'', or the Romans under Vespasian and Titus, who burned the
faithless city; and the psalm he cites for the angels is Ps~77, the Introit's
psalm, at its forty-ninth verse. Of the highways he writes, \latin{Gentilium
populus non erat in viis, sed in exitibus viarum}: the people of the Gentiles
was not in the ways, but at the ends of the ways; and he asks how any good
could be found among them, and answers with St Paul that the Gentiles who by
nature do the things of the law condemn the Jews (Rom 2:14; \work{In
Matthaeum} III, PL 26, cols.~159--160).

St Augustine answers the question of the highways in a short note of his
\work{Quaestiones evangeliorum}, and his answer takes a form of its own. The
wedding is the Word incarnate, \latin{quia in ipso homine suscepto Ecclesia Deo
copulata est}, because in the humanity he assumed the Church is joined to God;
the bulls are \latin{principes plebium}, the leaders of the peoples; and of the
command to go to the ends of the roads he writes, \latin{viæ intelliguntur
dogmata Gentium; quia ex omnibus illis ad nuptias venerunt, id est, Christo
crediderunt}: the ways are understood as the teachings of the Gentiles, because
out of all of them they came to the wedding, that is, they believed in Christ
(\work{Quaestiones evangeliorum} I.31; PL 35, col.~1329). For Augustine the
roads are not the peoples themselves but what they held; the call reaches men
in the doctrines they lived by and brings them out of all of them to Christ.
His answer points where Irenaeus's, Hilary's, Jerome's and Chrysostom's does,
to the Gentiles, and it says nothing of who were first invited or why they
refused, where the other four begin.

\subsection{The Introit's psalm and the Alleluia's command}

The Introit's verse opens the psalm in which Asaph tells Israel its history,
``That another generation might know them'' and ``may not become like their
fathers'' (Ps 77:6, 8). St Augustine reads the psalm's purpose as a warning to
those who hear it later: it contains ``the things which are said to have been
done among the old people: but the new and latter people is being admonished,
to beware that it be not ungrateful regarding the blessings of God''
(\work{Enarrationes} 77, 1). The Greek \work{Expositiones in Psalmos} printed
under the name of St Athanasius go further on the verse the Introit sings. They
ask to which people the command ``Attend, O my people'' is given, and answer
that it is the people gathered from the nations, and that the law it is to
attend to is the law of the Gospel. And the psalm's second verse, ``I will open
my mouth in parables'', is the verse St Matthew cites as fulfilled when Jesus
spoke in parables (Mt 13:35); the Gospel of this Mass opens with Jesus speaking
\latin{in parabolis} to the chief priests and Pharisees.

The Alleluia's verse ends with a command to declare God's deeds among the
Gentiles, and St Augustine hears in it the voice of prophecy addressing the
preachers of the Gospel. Having noted the other Latin copies'
\latin{Evangelizate}, preach the Gospel, he asks, ``Unto whom is this
addressed, save unto the Evangelists in prophecy?''
(\work{Enarrationes} 104, 1). Later in the same exposition he says who the
psalm's ``seed of Abraham'' are: \latin{non de sola gente Hebraeorum, sed
quotquot istam gratiam suscipiunt in omni terra}, not of the Hebrew nation
alone, but all who receive this grace in every land (104, 36). The Greek
\work{Expositiones} read the same verse as the Holy Spirit's command to the
apostles to make known to the people from the nations all the wonders the
Lord worked at his coming. Schuster names the apostles, and the bishops and
pastors of souls after them, as those who have always held this duty of
preaching the most important of all (p.~172), and St Robert Bellarmine reads the
verse as the psalmist's call to God's people to ``speak in all directions among
the gentiles of the wonderful works of God''.

Set before the Gospel, the Alleluia's command and the king's \latin{Ite ergo ad
exitus viarum} are one command in two forms, the psalm's prophecy and the
parable's story. St Augustine takes both steps, though in two separate works, reading the
Alleluia's verse of the evangelists in one and the Gospel's highways of the
Gentiles in the other. In this Mass the one follows the
other within a few lines, and the assembly hears the command before it hears
the story of its fulfilment.

The Alleluia's psalm also gives the reason for all the history it recites. Its
last verse is ``That they might observe his justifications, and seek after his
law'' (Ps 104:45), and St Augustine takes that verse as the key to the whole
psalm: the goods of the land were given so that God's servants might be free
for what is eternal, \latin{Ut custodiant justificationes ejus}
(\work{Enarrationes} 104, 34). The Communion of this Mass asks for the same
thing, that the ways of the one who prays be directed \latin{ad custodiendas
iustificationes tuas}. The history ends in keeping, and so does the Mass.

\subsection{The other elements on the way}

The Collect's prayer to be \latin{expediti}, disencumbered, in mind and body,
answers the excuse of the first invited. They went off \latin{alius in villam
suam, alius vero ad negotiationem suam}, and the witnesses name what held them.
For Hilary the farm is the ambition of the world and the merchandise the love
of money; for St Gregory the Great, \latin{In villam quippe ire est labori
terreno immoderate incumbere, in negotiationem vero ire est actionum
saecularium lucris inhiare}, to go to the farm is to lie immoderately upon
earthly labour, and to go to the merchandise is to gape after the profits of
worldly business (\work{Homiliae in Evangelia} 38, 5). The Collect asks to be
free of exactly that weight. The Epistle's ``for we are members one of
another'' is the rule of a body gathered from both peoples, and Jerome gives
the neighbour to whom truth is owed the widest possible extent:
\latin{omnem hominem, qui ex eodem nobiscum parente generatus est}, every man
born of the same first parent as we (on Eph 4:25, PL 26, col.~509).

The Gradual is the prayer of a people still on the road. St Augustine hears in
the psalm the voice of the whole Christ, head and members, which must go on
crying to God for as long as tribulation remains for the Church, even to the
end of the world (\work{Enarrationes} 140, 2). Hilary, on the Gradual's verse,
gives the evening sacrifice to us upon whom the end of the ages has come, the
people of the last age after the blood sacrifices have ceased (\work{Tractatus
super Psalmos} 140, 4). The Offertory's walk
``in the midst of tribulation'' is the road the servants took and the Church
still takes. Augustine hears the verse's confidence in the face of every enemy:
``Let mine enemies rage: what can they do? \dots\ For mine enemies cannot
separate me from Thee'' (\work{Enarrationes} 137, 11). The Secret brings the gifts of those who have
come in from the highways before the eyes of God's majesty, and the
Postcommunion asks that they cleave to the commandments \latin{semper}. The
call ends not in having been called, but in keeping what the caller commands.

\subsection{Two difficulties: the burned city, and St Gregory's highways}

The first difficulty is the burned city. St John Chrysostom reads it of
Jerusalem under Vespasian and Titus, and St Jerome keeps the Romans as one of
two possibilities beside the avenging angels of the Introit's psalm. St Gregory
reads it otherwise: the city burned is the flesh of the persecutors, which
suffers with their souls in the eternal fire, and the armies are the angels,
by whom every judgement is carried out; and he points his hearers to what
they can see for themselves, the vengeance that their fathers heard of and
that \latin{nos autem jam cernimus}, we now behold, since the proud persecutors
of the martyrs are forgotten while the martyrs' death flowers in the faith of
the living (\work{Homiliae} 38, 5). St Hilary, too, reads the armies as the armies
of heaven and the fire as eternal, and St Irenaeus says only that the king
burned their city, without naming it. The historical and the eschatological
readings stand side by side, and neither is the only answer. The parable's
hearers are those the Missal names, the chief priests and Pharisees who knew
he spoke of them (Mt 21:45). St Augustine sets the limit to any wider judgement
inside the psalm of Israel's history itself: ``there were even in that people
they that understood, having the faith which was afterwards revealed''
(\work{Enarrationes} 77, 2). And Chrysostom supplies the counterweight in his own homily, since he requires
the garment of the Gentiles who came in as strictly as he blames the first
invited for refusing.

The second difficulty is St Gregory the Great's reading of the highways, which gives the question a different answer. Gregory preached this
Gospel to the people of Rome in the basilica of St Clement, and in the whole
homily he names neither Israel nor the nations. His invited are \latin{amici},
friends, called first by the prophets and then by the apostles (38, 3); his
refusers are those who give themselves beyond measure to earthly labour and
worldly gain, and the worse of them persecute the grace of the one who calls
(38, 5). And the ends of the roads are not the Gentiles. \latin{Si in Scriptura
sacra vias actiones accipimus, exitus viarum intelligimus defectus actionum,
quia illi plerumque facile ad Deum veniunt, quos in terrenis actibus prospera
nulla comitantur}: if in Holy Scripture we take the ways as our doings, we
understand the ends of the ways as the failure of our doings, because those
whom no prosperity attends in their earthly affairs often come readily to God
(38, 6). This is a different answer, not a detail: for Gregory the highways
are the places where worldly undertakings run out, and the guests found there
are those whose affairs have failed. The two answers were already read side by
side in the ninth century. Rabanus Maurus's catena on St Matthew gives
Augustine's sentence on the teachings of the Gentiles first and Gregory's after
it, and heads Gregory's as the moral sense. Gregory's answer makes of the
highways a moral sense of the call, and it stands beside the answer of
Irenaeus, Hilary, Jerome and Chrysostom, for whom the highways are the nations,
and beside Augustine's own, for whom they are the teachings of the Gentiles.

\subsection{The four senses of this reading}

\begin{fourSenses}
\item[Literal.] In the temple, answering the chief priests and Pharisees who
knew that his parables were about them, Jesus tells of a king whose invited
guests refuse his son's wedding, maltreat and kill his servants, and see their
city burned, while the hall is filled from the highways with bad and good.
Ps~77 recites God's deeds so that the coming generation may not be like its
fathers, and Ps~104 bids Israel praise God and declare his deeds among the
Gentiles.

\item[Allegorical.] The king is the Father, and the wedding is the Son's union
with the Church through the Incarnation (Gregory; Augustine, for whom the
wedding is the Word incarnate, in whose assumed humanity the Church is joined
to God), a Church gathered from Jews and Gentiles (Jerome). The servants are
the prophets and then the apostles (Gregory; Chrysostom, who names John and the
Son between them), sent by the one God who called those of old by the prophets
and calls us by the apostles (Irenaeus); for Hilary they are the apostles and
then apostolic men, their successors. The highways are the nations (Irenaeus,
Hilary, Jerome, Chrysostom), or the teachings of the Gentiles, from all of
which those who believed came to Christ (Augustine). The Alleluia's command is
addressed to the evangelists in prophecy (Augustine) and to the apostles for
the people of the nations (the Greek \work{Expositiones}).

\item[Moral.] Not to make light of the call for farm or trade, which are
immoderate labour and worldly gain (Gregory) or ambition and the love of money
(Hilary), but to set the things of the spirit above even necessary things
(Chrysostom); not to be ungrateful as a people that has been told the story of
those before it (Augustine on Ps~77); not to trust in faith alone once inside
(Chrysostom); and, in Gregory's own form of the highways, to hear the call
when worldly undertakings fail, since those whom no prosperity attends often
come readily to God.

\item[Anagogical.] The wedding is the mystery of the eternal glory to be
received in the resurrection, and when all is ready and the multitude pleasing
to God is gathered into its number, the glory of the heavenly kingdom is
announced as a wedding (Hilary). The hall filled from the highways is the
Church of the nations awaiting that announcement. Of those gathered the few are
not few among the invited but rare among the chosen, \latin{non est \dots\
paucitas in invitatis, sed raritas in electis} (Hilary), as it was in the former
covenant (Irenaeus).
\end{fourSenses}
